\documentclass{beamer}

\usepackage{amsmath,amsfonts,amssymb}
\usepackage{physics}
\usepackage{bm}
\usepackage{quantikz}

\title{Reducing the Complexity of Unitary Coupled Cluster Circuits}
\subtitle{Tensor Factorizations and Parametric Matrix Models}
\author{}
\date{}

\begin{document}

\frame{\titlepage}

%================================================
\begin{frame}{Outline}

\begin{itemize}
\item Complexity of the UCC ansatz
\item Tensor structure of cluster amplitudes
\item Scaling of UCC circuits
\item Tensor factorization methods
\item Parametric matrix models
\item Circuit depth reduction
\end{itemize}

\end{frame}

%================================================
\section{Complexity of the UCC Ansatz}

%------------------------------------------------
\begin{frame}{Unitary Coupled Cluster}

The UCC wavefunction is

\[
|\Psi\rangle =
e^{T-T^\dagger}|\Phi_0\rangle
\]

\end{frame}

%------------------------------------------------
\begin{frame}{Cluster Operator}

For doubles

\[
T_2 =
\frac{1}{4}
\sum_{abij}
t_{ij}^{ab}
a_a^\dagger a_b^\dagger a_j a_i
\]

\end{frame}


%------------------------------------------------
\begin{frame}{Generator}

  Define

  \[
  \tau = T_2 - T_2^\dagger
  \]

  Thus

  \[
  U = e^\tau
  \]

\end{frame}

%------------------------------------------------
\begin{frame}{Number of Amplitudes}

  For $N$ orbitals

  \[
  t_{ij}^{ab}
  \]

  contains

  \[
  N^4
  \]

  elements.

\end{frame}

%------------------------------------------------
\begin{frame}{Implication}

  The number of excitation operators grows as

  \[
  N^4
  \]

  Therefore the number of circuit blocks grows rapidly.

\end{frame}

%================================================
\section{Tensor Structure of Amplitudes}

%------------------------------------------------
\begin{frame}{Amplitude Tensor}

  The cluster amplitudes form a rank-4 tensor

  \[
  t_{ij}^{ab}
  \]

  with indices

  \[
  i,j = \text{occupied}
  \]

  \[
  a,b = \text{virtual}
  \]

\end{frame}

%------------------------------------------------
\begin{frame}{Symmetry}

  The tensor satisfies

  \[
  t_{ij}^{ab} = -t_{ji}^{ab}
  \]

  \[
  t_{ij}^{ab} = -t_{ij}^{ba}
  \]

  This reduces the number of independent elements.

\end{frame}

%------------------------------------------------
\begin{frame}{Remaining Complexity}

  Despite antisymmetry

  \[
  N_{parameters} \sim N^4
  \]

\end{frame}

%================================================
\section{Circuit Implementation}

%------------------------------------------------
\begin{frame}{Jordan–Wigner Mapping}

  Fermionic operators are mapped to qubits

  \begin{align}
    a_p^\dagger &= \frac12(X_p-iY_p)\prod_{j<p}Z_j \\
    a_p &= \frac12(X_p+iY_p)\prod_{j<p}Z_j
  \end{align}

\end{frame}

%------------------------------------------------
\begin{frame}{Pauli Strings}

  Excitations become Pauli strings

  \[
  X_i X_j Y_k Y_l
  \]

\end{frame}

%------------------------------------------------
\begin{frame}{Circuit Block}

  Each excitation corresponds to

  \[
  e^{i\theta P}
  \]

  where $P$ is a Pauli string.

\end{frame}

%------------------------------------------------
\begin{frame}{Example Circuit}

  \[
  \begin{quantikz}
    \lstick{q_0} & \ctrl{1} & \qw & \qw & \qw \\
    \lstick{q_1} & \targ{} & \ctrl{1} & \qw & \qw \\
    \lstick{q_2} & \qw & \targ{} & \ctrl{1} & \qw \\
    \lstick{q_3} & \qw & \qw & \targ{} & \gate{R_z(\theta)}
  \end{quantikz}
  \]

\end{frame}

%================================================
\section{Tensor Factorization}

%------------------------------------------------
\begin{frame}{Low Rank Factorization}

  We approximate

  \[
  t_{ij}^{ab}
  \]

  as

  \[
  t_{ij}^{ab}
  =
  \sum_k
  X_{ia}^{(k)} X_{jb}^{(k)}
  \]

\end{frame}

%------------------------------------------------
\begin{frame}{Interpretation}

  Instead of $N^4$ parameters we use

  \[
  kN^2
  \]

  parameters.

\end{frame}

%------------------------------------------------
\begin{frame}{Reduced Operator}

  The cluster operator becomes

  \[
  T_2 =
  \sum_k
  \left(
  \sum_{ia} X_{ia}^{(k)} a_a^\dagger a_i
  \right)^2
  \]

\end{frame}

%================================================
\section{Operator Factorization}

%------------------------------------------------
\begin{frame}{Factorized Generator}

  Define

  \[
  B_k =
  \sum_{ia}
  X_{ia}^{(k)} a_a^\dagger a_i
  \]

  Then

  \[
  T_2 =
  \sum_k B_k^2
  \]

\end{frame}

%------------------------------------------------
\begin{frame}{Unitary Operator}

  \[
  U =
  e^{\sum_k (B_k^2-B_k^{\dagger 2})}
  \]

\end{frame}

%================================================
\section{Parametric Matrix Models}

%------------------------------------------------
\begin{frame}{PMM Ansatz}

  Parametric matrix models write

  \[
  T(\theta)=
  \sum_k \theta_k M_k
  \]

\end{frame}

%------------------------------------------------
\begin{frame}{Reduced Generator}

  \[
  \tau =
  \sum_k \theta_k (M_k-M_k^\dagger)
  \]

\end{frame}

%------------------------------------------------
\begin{frame}{Circuit Representation}

  \[
  U(\theta)=
  \prod_k
  e^{\theta_k(M_k-M_k^\dagger)}
  \]

\end{frame}

%================================================
\section{Circuit Complexity}

%------------------------------------------------
\begin{frame}{Original UCC}

  Parameter count

  \[
  N^4
  \]

  Circuit depth

  \[
  D \sim \frac{N^4}{\Delta t}
  \]

\end{frame}

%------------------------------------------------
\begin{frame}{PMM}

  Parameter count

  \[
  N^2
  \]

  Circuit depth

  \[
  D \sim \frac{N^2}{\Delta t}
  \]

\end{frame}

%------------------------------------------------
\begin{frame}{Improvement}

  Reduction factor

  \[
  N^2
  \]

\end{frame}

%================================================
\section{Physical Interpretation}

%------------------------------------------------
\begin{frame}{Collective Excitations}

  PMM operators correspond to

  collective particle-hole excitations.

\end{frame}

%------------------------------------------------
\begin{frame}{Relation to Many-Body Theory}

  The approach resembles

  \begin{itemize}
  \item random phase approximation
  \item low-rank factorization
  \item tensor decomposition
  \end{itemize}

\end{frame}

%================================================
\section{Summary}

%------------------------------------------------
\begin{frame}{Key Results}

  UCC circuits become large because

  \[
  t_{ij}^{ab}
  \]

  contains

  \[
  N^4
  \]

  parameters.

\end{frame}

%------------------------------------------------
\begin{frame}{Solutions}

  Circuit depth can be reduced using

  \begin{itemize}
  \item tensor factorization
  \item adaptive ansätze
  \item parametric matrix models
  \end{itemize}

\end{frame}

%------------------------------------------------
\begin{frame}{Final Scaling}

  \[
  D_{UCC} \sim \frac{N^4}{\Delta t}
  \]

  \[
  D_{PMM} \sim \frac{N^2}{\Delta t}
  \]

\end{frame}

\end{document}
